\documentclass[a4paper]{ctexart}

% 导言区，加载宏包和各项设置
\usepackage[hmargin=1.25in,vmargin=1in]{geometry}  %  word 版式
\usepackage[x11names]{xcolor} %must before tikz, x11names defines RoyalBlue3
\usepackage[bookmarksnumbered,unicode, pdfborder=1,breaklinks,colorlinks,linkcolor=RoyalBlue3,urlcolor=blue]{hyperref}
\usepackage{graphicx} % 图片
\usepackage{fancyhdr} % 页眉页脚
\usepackage[raggedright]{titlesec} % 标题
%\usepackage{enumitem}
\usepackage{listings}
\usepackage{tikz}

% 环境设置
\graphicspath{{figures/}} % 设置图片路径

% 页眉页脚格式
\pagestyle{fancy}
\fancyhf{}
\cfoot{\thepage}
\rhead{USB~启动盘的制作}
\lhead{\includegraphics[width=3.5cm]{logo/cs2c-logo}}
\renewcommand{\headrulewidth}{0.4pt}

% 行距
\linespread{1.6}

\lstset{
basicstyle=\ttfamily,
columns=flexible,
numbers=left,
numberstyle=\footnotesize,
keywordstyle=\color{blue!70},
commentstyle=\color{red!50!green!50!blue!50},
frame=single,backgroundcolor=\color{white}, 
framexleftmargin=2em,
breaklines=true,
escapeinside=``,
xleftmargin=3.5em,
xrightmargin=1em, 
aboveskip=1em
}

% 正文
\begin{document}

%%%%　cover %%%%
\thispagestyle{empty}

\noindent\begin{minipage}{\textwidth}
	\begin{flushleft}
	\includegraphics{logo/cs2c-short-logo}
	\end{flushleft}
\end{minipage}

\vspace{\stretch{1}}
\noindent\begin{minipage}{\textwidth}
\centering
{\LARGE \bfseries USB~启动盘的制作}
\noindent\rule[1.5ex]{\textwidth}{1pt}
\end{minipage}

\vspace{\stretch{2}}
\noindent\begin{center}
2016~年~07~月~14~日
\end{center}
\clearpage

%%%% main body %%%%

\section{需求说明}
Linux~各个发行版本是通过~LiveCD ISO~系统镜像方式发布。对于频繁测试系统镜像的工作而言，如果通过刻录光盘方式则需要消耗大量时间及光盘。本文将介绍一种通过~USB（3.0~接口）启动盘的方式实现系统的快速安装。

\section{软硬件要求}
%\begin{enumerate}[itemsep=0pt,parsep=0pt,label={\chinese*、}]
%\item \textbf{软件：} Linux~系统，Grub
%\item \textbf{硬件：}  ~3.0接口的~U~盘\footnote{usb~2.0~接口的也可以，但是其传输速度较低，不建议使用}。
%\end{enumerate}
\textbf{软件：}~Linux~系统、Grub；

\textbf{硬件：}~3.0接口的~U~盘\footnote{usb~2.0~接口的也可以，但是其传输速度较低，不建议使用}。

\section{制作步骤}
下面是以~U~盘（后面约定设备显示~sdb~）并使用第一个分区sdb1作为存储启动文件分区为例的详细制作步骤。

\subsection{U~盘挂载及~grub~安装}

\begin{lstlisting}[language=bash]
mount /dev/sdb1 /mnt/usb
grub2-install --force --no-floppy --root-directory=/mnt/usb /dev/sdb  #有些系统可能为 grub-install
\end{lstlisting}

\emph{注意：/mnt/usb~为挂载点，大部分~U~盘会自动挂载在~/run/media/~下，可根据情况调整。}

命令执行完成后，则会在挂载点根下生成~boot~目录，同时手动在挂载点根下创建iso目录。将对应的镜像文件拷贝到该目录下，镜像名称需要和~grub.cfg~文件中写入的文件名保持一致。

目录结构：

\begin{figure}[htbp]
\setlength{\leftskip}{3em}
\begin{tikzpicture}
\node (boot) at (1,0) {boot};
\node (grub2) at (2,-0.5) {grub2};
\node (iso) at (1, -1) {iso};
\draw (0,0) -- (boot) -- (1, -0.5) -- (grub2);
\draw (0,0.2) -- (0,-1.2);
\draw (0, -1) -- (iso);
\end{tikzpicture}
\end{figure}


\subsection{grub.cfg~文件调整}
拷贝一个~grub.cfg~文件到~U~盘~boot/grub2~目录下，或则使用\verb|grub2-mkconfig|临时生成一个~grub.cfg~文件，添加或修改~menuentry~如下即可。

\begin{lstlisting}
menuentry 'TEST' {
    load_video
    set gfxpayload=keep
    insmod gzio
    insmod part_msdos
    insmod ext2
    set root='hd0,msdos1'
    loopback loop (hd0,1)/iso/test.iso
    linux16 (loop)/isolinux/vmlinuz0 root=live:CDLABEL=NeoKylin-Linux-Desktop-XX rootfstype=auto ro rd.live.image iso-scan/filename=/iso/test.iso quiet rhgb
    initrd16 (loop)/isolinux/initrd0.img
}
\end{lstlisting}

\emph{注意：root=‘hd0,msdos*’~和~(hd0,*)~需要和使用的分区对应；CDLABEL=*~需要和镜像文件对应；镜像名称~test.iso~可自定义，需要和下一步骤中的~iso~目录下的文件对应。}


\end{document}